% ============================================================
% Bölüm 9: INNATE Mimarisi
% ============================================================

\chapter{INNATE Mimarisi}
\label{ch:innate_mimari}

INNATE (Interpretable Neural Network Architecture for Turbulence Emulation), klasik neural operator'lardan temel bir farkla ayrılır: ağdaki her nöron bir \textit{fiziksel operatör}dür. MLP katmanları, attention mekanizmaları veya kara kutu bileşenler yoktur. \"Oğrenilebilir parametreler doğrudan fiziksel anlamlara sahiptir: etkin viskozite, adveksiyon modulatörü, kaldırma kuvveti katsayısı gibi.


% ============================================================
\section{Temel Prensip: Fizik = Ağ Yapısı}
\label{sec:temel_prensip}

Geleneksel neural network'lerde:
\[
    h = \sigma(Wx + b)
\]
nöronu, öğrenilen ağırlık matrisi $W$ ve bias $b$ ile girdiye lineer dönüşüm + aktivasyon uygular. Bu nöron \textit{evrensel} ama \textit{yorumlanamazdır}: $W$'nun fiziksel anlamı yoktur.

INNATE'te her nöron, NS denklemlerinin belirli bir terimini hesaplar:
\begin{align}
    \text{Advection nöronu:} &\quad -\alpha \cdot (\vect{u} \cdot \gradient)\vect{u} \\
    \text{Diffusion nöronu:} &\quad \nu_{\text{eff}} \cdot \laplacian \vect{u} \\
    \text{Projection nöronu:} &\quad \vect{u} - \gradient \phi \quad (\divergence \vect{u} = 0) \\
    \text{Buoyancy nöronu:} &\quad \beta_{\text{eff}} \cdot \Rich \cdot T' \cdot \hat{\vect{e}}_y
\end{align}
Her nöronün bir veya birkaç öğrenilebilir parametresi vardır ($\alpha$, $\nu_{\text{eff}}$, $\beta_{\text{eff}}$, \ldots). Eğitim sonunda bu parametreler fiziksel olarak yorumlanabilir: ``model hangi etkin viskoziteyi öğrendi?'' sorusu doğrudan cevaplanabilir.

\begin{remark}
Parametre sayısı perspektifi: INNATE yaklaşık $\sim 10$K parametreye sahipken, tipik bir FNO modeli $\sim 10^6$ parametreye sahiptir. Bu $100\times$ fark, INNATE'in yorumlanabilirliğinin ve fiziksel çerçevesinin doğrudan sonucudur.
\end{remark}


% ============================================================
\section{Nöron Tipleri}
\label{sec:noron_tipleri}

Aşağıda her nöron tipi, matematiksel formu, fiziksel anlamı ve öğrenilebilir parametreleri ile birlikte tanımlanmıştır.


% --- 1. Advection3D ---
\subsection{Advection3D: Nonlineer Konveksiyon}
\label{sec:advection_noron}

\textbf{Fiziksel anlam:} Akışkanın momentumunu kendi hız alanıyla taşıması.

\textbf{Matematiksel form (convective form):}
\begin{equation}
    \mathcal{A}(\vect{u}) = -\alpha_{\text{adv}} \cdot (\vect{u} \cdot \gradient)\vect{u}
    \label{eq:advection_neuron}
\end{equation}
Bileşen yazımı:
\begin{equation}
    \mathcal{A}_i = -\alpha_{\text{adv}} \cdot \sum_{j=1}^{3} u_j \pd{u_i}{x_j}, \quad i \in \{1, 2, 3\}
    \label{eq:advection_component}
\end{equation}

\textbf{Hesaplama adımları:}
\begin{enumerate}
    \item Türevler spektral olarak hesaplanır: $\widehat{\partial u_i / \partial x_j} = i k_j \hat{u}_i$.
    \item Fiziksel uzaya döndürülür: $\partial u_i / \partial x_j = \text{IFFT}(\ldots)$.
    \item Çarpım fiziksel uzayda yapılır: $u_j \cdot (\partial u_i / \partial x_j)$.
    \item \textbf{Dealiasing} uygulanır (2/3 kuralı veya mode-index bazlı).
\end{enumerate}

\textbf{Dealiasing neden gerekli?} Fiziksel uzayda iki alanın çarpımı, Fourier uzayında konvolüsyona karşılık gelir. Eğer her iki alan $k_{\max}$'a kadar modlara sahipse, çarpım $2k_{\max}$'a kadar modlar üretir. Grid sadece $k_{\max}$'a kadar çözebildiğinden, $k_{\max}$'ı aşan modlar düşük modlara ``katlanır'' (aliasing). 2/3 kuralı, $k > (2/3)k_{\max}$ modlarını sıfırlayarak bu hatanın birikmesini önler.

\textbf{Öğrenilebilir parametre:}
\begin{itemize}
    \item \texttt{advection\_modulator} $(\alpha_{\text{adv}})$: başlangıç değeri $= 1.0$, klamplama: $[0.5, 2.0]$.
\end{itemize}
$\alpha_{\text{adv}} = 1$ tam adveksiyon demektir. Model bunu biraz büyütebilir veya küçültebilir: LES gridinde çözülemeyen ince yapıların adveksiyon üzerindeki etkisini modelle.


% --- 2. Diffusion ---
\subsection{Diffusion: Visköz Yayılma}
\label{sec:diffusion_noron}

\textbf{Fiziksel anlam:} Moleküler viskozite tarafından momentumun yayılması.

\textbf{Matematiksel form:}
\begin{equation}
    \mathcal{D}(\vect{u}) = \frac{\nu_{\text{scale}}}{\Rey} \laplacian \vect{u}
    \label{eq:diffusion_neuron}
\end{equation}

\textbf{Hesaplama:} Tamamen spektral uzayda:
\begin{equation}
    \widehat{\laplacian u_i} = -(k_x^2 + k_y^2 + k_z^2) \hat{u}_i = -|\vect{k}|^2 \hat{u}_i
    \label{eq:spectral_laplacian}
\end{equation}
Bu işlem FFT + çarpım + IFFT olarak çok ucuzdur.

\textbf{Öğrenilebilir parametre:}
\begin{itemize}
    \item \texttt{diffusion\_scale} $(\nu_{\text{scale}})$: başlangıç $= 1.0$, klamplama: $[0.1, 3.0]$.
\end{itemize}
$\nu_{\text{scale}} = 1$ dolaylı moleküler viskoziteyi verir. Model bunu $> 1$ yaparak etkin viskoziteyi artırabilir --- LES grid'inde çözülemeyen küçük ölçeklerin dissipasyonunu kaba bir şekilde modellemek için.


% --- 3. EddyViscosity3D ---
\subsection{EddyViscosity3D: LES Nöronu}
\label{sec:eddy_noron}

\textbf{Fiziksel anlam:} Çözülemeyen küçük ölçeklerin büyük ölçekler üzerindeki etkisini modeller (subgrid-scale modeli).

\textbf{Matematiksel form (Smagorinsky):}
\begin{equation}
    \mathcal{E}(\vect{u}) = \divergence \left( \nu_t \, \gradient \vect{u} \right)
    \label{eq:eddy_neuron}
\end{equation}
\begin{equation}
    \nu_t = (C_s \, \bar{\Delta})^2 \, |\tensor{S}|
    \label{eq:smagorinsky}
\end{equation}
\begin{itemize}
    \item $C_s$: Smagorinsky sabiti (öğrenilebilir)
    \item $\bar{\Delta} = (\Delta x \cdot \Delta y \cdot \Delta z)^{1/3}$: filtre genişliği (grid boyutlarından hesaplanır)
    \item $|\tensor{S}| = \sqrt{2 S_{ij} S_{ij}}$: strain rate büyüklüğü
    \item $S_{ij} = \frac{1}{2}\left(\pd{u_i}{x_j} + \pd{u_j}{x_i}\right)$: strain rate tensörü
\end{itemize}

\textbf{Neden kritik?} $96 \times 160 \times 64$ gridi, $\Rey = 10\,000$'de tüm ölçekleri çözecek incelikte değildir. Grid analizi LES skorunu $74/100$ olarak belirlemiştir: makul ama DNS değil. EddyViscosity kapalıyken, çözülemeyen modlara enerji birikir ve 63 saatlik eğitimde görüldüğü gibi patlama (NaN) meydana gelir.

\textbf{Öğrenilebilir parametre:}
\begin{itemize}
    \item \texttt{cs\_value} $(C_s)$: başlangıç $= 0.15$, klamplama: $[0.05, 0.30]$.
\end{itemize}
Klasik Smagorinsky'de $C_s \approx 0.17$ (izotropik türbülans için). INNATE'in bunu öğrenmesine izin verilir: mixed convection'da kanal akışı ve plume yapıları birlikte bulunduğundan, optimal $C_s$ homojen türbülanstan sapabilir.


% --- 4. Projection3D ---
\subsection{Projection3D: Süreklilik Zorlaması}
\label{sec:projection_noron}

\textbf{Fiziksel anlam:} Sıkıştırılamaz akışta $\divergence \vect{u} = 0$ koşulunu zorlar.

\textbf{Matematiksel form (Helmholtz ayrıştırması):}

Herhangi bir hız alanı $\vect{u}^*$ (henüz divergence-free olmayan), şöyle ayrıştırılır:
\begin{equation}
    \vect{u}^* = \vect{u} + \gradient \phi
    \label{eq:helmholtz}
\end{equation}
Burada $\vect{u}$ solenodial ($\divergence \vect{u} = 0$) ve $\gradient \phi$ irrotasyonel kısımdır. Sürekliliği zorlamak için:
\begin{enumerate}
    \item Poisson denklemi çözülür:
    \begin{equation}
        \laplacian \phi = \divergence \vect{u}^*
        \label{eq:poisson}
    \end{equation}
    \item Hız düzeltilir:
    \begin{equation}
        \vect{u} = \vect{u}^* - s_p \cdot \gradient \phi
        \label{eq:projection_correction}
    \end{equation}
\end{enumerate}

\textbf{Spektral çözüm:} Periyodik domain'de Poisson denklemi Fourier uzayında tek bir bölme işlemine dönüşür:
\begin{equation}
    \hat{\phi}(\vect{k}) = \frac{\widehat{\divergence \vect{u}^*}(\vect{k})}{-|\vect{k}|^2}, \quad |\vect{k}| \neq 0; \qquad \hat{\phi}(\vect{0}) = 0
    \label{eq:poisson_spectral}
\end{equation}
Sıfır modu ($\vect{k} = \vect{0}$) ortalama basınca karşılık gelir ve periyodik domain'de belirsizdir; sıfıra sabitlenir.

\textbf{Öğrenilebilir parametre:}
\begin{itemize}
    \item \texttt{pressure\_scale} $(s_p)$: başlangıç $= 1.0$, klamplama: $[0.8, 1.2]$.
\end{itemize}
$s_p = 1$ tam Projection demektir. Model bunu hafifçe ayarlayarak, unrolled time-stepping'de basınç düzeltmesinin ``şiddetini'' optimize edebilir.


% --- 5. Forcing3D ---
\subsection{Forcing3D: Dış Kuvvet}
\label{sec:forcing_noron}

\textbf{Fiziksel anlam:} Kolmogorov body forcing ile rüzgâr etkisi.

\textbf{Matematiksel form:}
\begin{equation}
    \mathcal{F}(y) = A \sin\!\left(\frac{k_f \cdot 2\pi y}{L_y}\right) \hat{\vect{e}}_x
    \label{eq:forcing_neuron}
\end{equation}

\textbf{Öğrenilebilir parametre:}
\begin{itemize}
    \item \texttt{amplitude} $(A)$: başlangıç $= 1.0$, klamplama: $[0.01, 10.0]$.
\end{itemize}
$A$'nın büyüklüğü, enerji enjeksiyon oranını ($\varepsilon_{\text{in}} = \langle \vect{u} \cdot \vect{F} \rangle$) kontrol eder. Eğitim sırasında model, enerji dengesini sağlayacak $A$'yı öğrenir.


% --- 6. Buoyancy3D ---
\subsection{Buoyancy3D: Kaldırma Kuvveti}
\label{sec:buoyancy_noron}

\textbf{Fiziksel anlam:} Sıcaklık pertürbasyonunun dikey momentuma etkisi.

\textbf{Matematiksel form:}
\begin{equation}
    \mathcal{B}(T') = \beta_{\text{eff}} \cdot \Rich \cdot T' \cdot \hat{\vect{e}}_y
    \label{eq:buoyancy_neuron}
\end{equation}

Yalnızca $v$-denklemine katılır. $T' > 0$ (ortalamadan sıcak) olan bölgelerde yukarı kuvvet, $T' < 0$ olan bölgelerde aşağı kuvvet oluşturur.

\textbf{Öğrenilebilir parametre:}
\begin{itemize}
    \item \texttt{buoyancy\_strength} $(\beta_{\text{eff}})$: başlangıç $= 1.0$, klamplama: $[0.0, 50.0]$.
\end{itemize}


% --- 7. ThermalAdvection3D ---
\subsection{ThermalAdvection3D: Sıcaklık Adveksiyonu}
\label{sec:thermal_adv_noron}

\textbf{Fiziksel anlam:} Hız alanının sıcaklık pertürbasyonunu taşıması.

\textbf{Matematiksel form:}
\begin{equation}
    \mathcal{T}_{\text{adv}}(\vect{u}, T') = -(\vect{u} \cdot \gradient)T' = -\left(u\pd{T'}{x} + v\pd{T'}{y} + w\pd{T'}{z}\right)
    \label{eq:thermal_adv_neuron}
\end{equation}

Hesaplama, Advection3D'nin skaler versiyonudur: türevler spektral, çarpımlar fiziksel uzayda, dealiasing uygulanır.


% --- 8. ThermalDiffusion3D ---
\subsection{ThermalDiffusion3D: Sıcaklık Difüzyonu}
\label{sec:thermal_diff_noron}

\textbf{Fiziksel anlam:} Moleküler iletimle sıcaklığın yayılması.

\textbf{Matematiksel form:}
\begin{equation}
    \mathcal{T}_{\text{diff}}(T') = \frac{\kappa_{\text{scale}}}{\Rey \, \Pran} \laplacian T'
    \label{eq:thermal_diff_neuron}
\end{equation}

\textbf{Öğrenilebilir parametre:}
\begin{itemize}
    \item \texttt{kappa\_scale} $(\kappa_{\text{scale}})$: başlangıç $= 1.0$, klamplama: $[0.1, 5.0]$.
\end{itemize}
$\kappa_{\text{scale}} > 1$ etkin termal difuziviteyi artırır, yani çözülemeyen küçük ölçekli termal yapıların difuzif etkisini modeller.


% ============================================================
\section{Parametre Tablosu}
\label{sec:parametre_tablosu}

\begin{table}[H]
\centering
\caption{INNATE nöron parametreleri özeti.}
\label{tab:neuron_params}
\begin{tabular}{@{}llccc@{}}
\toprule
\textbf{Nöron} & \textbf{Parametre} & \textbf{Init} & \textbf{Min} & \textbf{Max} \\
\midrule
Advection3D        & \texttt{advection\_modulator}  & $1.0$  & $0.5$  & $2.0$ \\
Diffusion          & \texttt{diffusion\_scale}      & $1.0$  & $0.1$  & $3.0$ \\
EddyViscosity3D    & \texttt{cs\_value}             & $0.15$ & $0.05$ & $0.30$ \\
Projection3D       & \texttt{pressure\_scale}       & $1.0$  & $0.8$  & $1.2$ \\
Forcing3D          & \texttt{amplitude}             & $1.0$  & $0.01$ & $10.0$ \\
Buoyancy3D         & \texttt{buoyancy\_strength}    & $1.0$  & $0.0$  & $50.0$ \\
ThermalAdvection3D & (parametresiz)                 & ---    & ---    & --- \\
ThermalDiffusion3D & \texttt{kappa\_scale}          & $1.0$  & $0.1$  & $5.0$ \\
\bottomrule
\end{tabular}
\end{table}

\begin{remark}
Başlangıç değerleri, nöronların \textit{tam fizik} (\textit{yani DNS benzeri}) davranışına karşılık gelecek şekilde seçilmiştir. $\alpha_{\text{adv}} = 1$, $\nu_{\text{scale}} = 1$, $\beta_{\text{eff}} = 1$ demek, denklemlerin hiçbir modifikasyon olmadan çözülmesi anlamına gelir. Eğitim, bu değerleri grid çözünürlüğüne ve LES etkilerine göre ince ayarlar.
\end{remark}


% ============================================================
\section{SpectralOps3DAniso: Spektral Operatörler}
\label{sec:spectral_ops}

Tüm nöronların türev hesapları, ortak bir spektral operatörler sınıfı üzerinden yapılır: \texttt{SpectralOps3DAniso}.

\subsection{Neden Anizotropik?}

Domain boyutları eşit değildir ($6 \times 10 \times 4$). Dolayısıyla her yöndeki dalga sayıları farklıdır:
\begin{equation}
    k_x = \frac{2\pi n_x}{L_x}, \quad
    k_y = \frac{2\pi n_y}{L_y}, \quad
    k_z = \frac{2\pi n_z}{L_z}
    \label{eq:aniso_wavenumbers}
\end{equation}
$n_x \in [0, N_x/2]$, $n_y \in [0, N_y/2]$, $n_z \in [0, N_z/2]$ tamsayı mod indeksleridir. Farklı $L_i$ değerleri, aynı mod indeksinin farklı fiziksel dalga sayılarına karşılık gelmesine neden olur.

\subsection{Temel Operatörler}

\begin{description}
    \item[\texttt{gradient($\hat{f}$)}:] Skaler alanın gradyanı:
    $\widehat{\partial f / \partial x_j} = i k_j \hat{f}$.

    \item[\texttt{laplacian($\hat{f}$)}:] Skaler/vektörel Laplacian:
    $\widehat{\laplacian f} = -(k_x^2 + k_y^2 + k_z^2) \hat{f}$.

    \item[\texttt{curl($\hat{\vect{u}}$)}:] Vortisite:
    $\hat{\omega}_i = i \epsilon_{ijk} k_j \hat{u}_k$ (Levi-Civita sembolü ile).

    \item[\texttt{divergence($\hat{\vect{u}}$)}:] $\widehat{\divergence \vect{u}} = i(k_x \hat{u} + k_y \hat{v} + k_z \hat{w})$.

    \item[\texttt{dealias($\hat{f}$)}:] Yüksek modları sıfırla. Mode-index bazlı: $|n_x| > N_x/3$, $|n_y| > N_y/3$ veya $|n_z| > N_z/3$ olan modlar kaldırılır.

    \item[\texttt{solve\_poisson($\hat{f}$)}:] Poisson denklemi:
    $\hat{\phi} = \hat{f} / (-|\vect{k}|^2)$, $\hat{\phi}(\vect{0}) = 0$.
\end{description}

\subsection{Dealiasing: Mode-Index ve Fiziksel $k$}

\"Onemli bir ayrım: Dealiasing işleminde \textit{mode-index} bazlı filtreleme kullanılır, \textit{fiziksel dalga sayısı büyüklüğü} ($|\vect{k}|$) bazlı değil. Nedeni: anizotropik domain'de aynı fiziksel $k$ değeri, farklı yönlerde farklı mod indekslerine karşılık gelir. Mode-index bazlı kural, FFT'nin ürettiği dizin yapısıyla doğrudan uyumludur.


% ============================================================
\section{Forward Pass Akışı}
\label{sec:forward_pass}

Bir zaman adımının INNATE içindeki akışı:

\begin{figure}[H]
\centering
\includegraphics[width=0.9\textwidth]{innate_architecture.png}
\caption{INNATE forward pass akış diyagramı. Her kutu bir fizik nöronunu, yanındaki sayılar öğrenilebilir parametre sayısını gösterir. Akış: momentum $\to$ termal $\to$ projeksiyon $\to$ zaman entegrasyonu.}
\label{fig:innate_architecture}
\end{figure}

\subsection{State Yapısı}

Giriş ve çıkış: \texttt{FluidState3D} --- hız bileşenleri ($u, v, w$), basınç ($p$) ve sıcaklık pertürbasyonu ($T'$).

\subsection{Bir Zaman Adımının Akışı}

\begin{enumerate}
    \item \textbf{Momentum RHS hesabı:}
    \begin{enumerate}
        \item Advection3D: $\vect{R}_{\text{adv}} = -\alpha_{\text{adv}} \cdot (\vect{u} \cdot \gradient)\vect{u}$
        \item Diffusion: $\vect{R}_{\text{diff}} = (\nu_{\text{scale}} / \Rey) \, \laplacian \vect{u}$
        \item EddyViscosity3D: $\vect{R}_{\text{eddy}} = \divergence(\nu_t \gradient \vect{u})$
        \item Forcing3D: $\vect{R}_{\text{force}} = A \sin(k_f \cdot 2\pi y / L_y) \, \hat{\vect{e}}_x$
        \item Buoyancy3D: $\vect{R}_{\text{buoy}} = \beta_{\text{eff}} \cdot \Rich \cdot T' \, \hat{\vect{e}}_y$
    \end{enumerate}
    Toplam: $\vect{R}_u = \vect{R}_{\text{adv}} + \vect{R}_{\text{diff}} + \vect{R}_{\text{eddy}} + \vect{R}_{\text{force}} + \vect{R}_{\text{buoy}}$

    \item \textbf{Sıcaklık RHS hesabı:}
    \begin{enumerate}
        \item ThermalAdvection3D: $R_{T,\text{adv}} = -(\vect{u} \cdot \gradient)T'$
        \item ThermalDiffusion3D: $R_{T,\text{diff}} = (\kappa_{\text{scale}} / \Rey\Pran) \, \laplacian T'$
        \item Kaynak terim: $R_{T,\text{src}} = v \cdot \Delta T / L_y$
    \end{enumerate}
    Toplam: $R_T = R_{T,\text{adv}} + R_{T,\text{diff}} + R_{T,\text{src}}$

    \item \textbf{Zaman entegrasyonu (TimeMarcher3D):}
    RK4 (4. derece Runge--Kutta) şeması:
    \begin{align}
        \vect{k}_1 &= \vect{R}(\vect{q}^n) \\
        \vect{k}_2 &= \vect{R}\!\left(\vect{q}^n + \tfrac{\Delta t}{2} \vect{k}_1\right) \\
        \vect{k}_3 &= \vect{R}\!\left(\vect{q}^n + \tfrac{\Delta t}{2} \vect{k}_2\right) \\
        \vect{k}_4 &= \vect{R}(\vect{q}^n + \Delta t \, \vect{k}_3) \\
        \vect{q}^{n+1} &= \vect{q}^n + \frac{\Delta t}{6}(\vect{k}_1 + 2\vect{k}_2 + 2\vect{k}_3 + \vect{k}_4)
    \end{align}
    Burada $\vect{q} = (u, v, w, T')$ tüm state değişkenlerini temsil eder.

    \item \textbf{Projection (div=0 zorlaması):}
    RK4 sonrasında elde edilen $\vect{u}^*$'a Projection3D uygulanır:
    \[
        \laplacian \phi = \divergence \vect{u}^*, \quad
        \vect{u}^{n+1} = \vect{u}^* - s_p \cdot \gradient \phi
    \]

    \item \textbf{Gauge fix:} $T' \leftarrow T' - \langle T' \rangle$.

    \item \textbf{Çıkış:} \texttt{FluidState3D}$(u^{n+1}, v^{n+1}, w^{n+1}, p^{n+1}, T'^{n+1})$.
\end{enumerate}

\begin{remark}
1 zaman adımı = 1 ağ katmanı. 20 adımlık unrolled eğitim $=$ 20 katmanlı bir ağdan geçiş. Geri yayılım (backpropagation), bu 20 katmandan geriye doğru gradient hesaplar.
\end{remark}


% ============================================================
\section{Yol C: Lokal-Adaptif Parametreler}
\label{sec:yol_c}

Temel INNATE'te her nöron \textit{tek global skaler} parametre kullanır: $\alpha_{\text{adv}}$ her yerde aynıdır. Bu, türbülant akışta yetersiz kalabilir, çünkü lokal koşullar (plume bölgesi vs.\ shear bölgesi) çok farklıdır.

\subsection{Lokal \"{I}nvariantlar}

Her grid noktasında, akış alanından doğrudan hesaplanabilen altı skaler büyüklük:
\begin{equation}
\begin{aligned}
    I_1 &= |\tensor{S}| = \sqrt{2 S_{ij} S_{ij}} && \text{(strain rate büyüklüğü)} \\
    I_2 &= |\vect{\omega}| = \sqrt{\omega_i \omega_i} && \text{(vortisite büyüklüğü)} \\
    I_3 &= Q = \frac{|\vect{\omega}|^2 - |\tensor{S}|^2}{2} && \text{($Q$-kriteri: girdap tanımlama)} \\
    I_4 &= |\gradient T'| && \text{(sıcaklık gradyanı)} \\
    I_5 &= E_{\text{local}} = \tfrac{1}{2}(u^2 + v^2 + w^2) && \text{(lokal kinetik enerji)} \\
    I_6 &= T' && \text{(sıcaklık pertürbasyonu)}
\end{aligned}
\label{eq:invariants}
\end{equation}

Bu büyüklükler, FFT türevlerinden hesaplanır ve otomatik türev (backprop) ile gradient akışını destekler.

\subsection{Adaptif Parametre Formülü}

Global skaler yerine, parametreler her noktada invariantlardan hesaplanır:

\textbf{Lineer versiyon:}
\begin{equation}
    \alpha_{\text{eff}}(\vect{x}) = \alpha_0 + \sum_{m=1}^{6} \alpha_m \, \hat{I}_m(\vect{x})
    \label{eq:linear_modulation}
\end{equation}
7 öğrenilebilir parametre per operatör per katman. $\hat{I}_m$ normalize edilmiş invariantlardır:
\begin{equation}
    \hat{I}_m = \frac{I_m}{\max(I_{m,\text{global}}, \, \varepsilon)}
    \label{eq:invariant_normalization}
\end{equation}
Normalizasyon olmadan invariantlar çok farklı ölçeklerdedir ($|S| \sim \mathcal{O}(100)$, $T' \sim \mathcal{O}(0.01)$) ve eğitim kararsızlaşır.

\textbf{Kuadratik versiyon:}
\begin{equation}
    \alpha_{\text{eff}} = \alpha_0
    + \sum_{m=1}^{6} \alpha_m \hat{I}_m
    + \sum_{m=1}^{3} \alpha_{6+m} \hat{I}_m^2
    + \alpha_{10} \hat{I}_1 \hat{I}_2
    + \alpha_{11} \hat{I}_1 \hat{I}_3
    + \alpha_{12} \hat{I}_2 \hat{I}_3
    \label{eq:quadratic_modulation}
\end{equation}
13 parametre per operatör per katman. Çapraz terimler ($\hat{I}_1 \hat{I}_2$ gibi), strain-vorticity etkileşimini yakalar.

\textbf{Pozitiflik zorlaması (difüziviteler için):}
\begin{equation}
    \nu_{\text{eff}}(\vect{x}) = \text{softplus}\!\left(\nu_0 + \sum_m \nu_m \hat{I}_m + \ldots\right)
    \label{eq:softplus_constraint}
\end{equation}
Softplus: $\text{softplus}(x) = \ln(1 + e^x)$, her zaman $> 0$.

\subsection{Blok Yapısı}

Yol C mimarisinde, 20--25 katman 3--4 bloğa gruplandırılır. Her blokta:
\begin{enumerate}
    \item FFT ile türevler ve invariantlar hesaplanır (blok başında bir kez).
    \item 5--7 katmanda RHS hesabı yapılır (sadece pointwise çarpım ve toplam; FFT yok).
    \item Projection ile $\divergence \vect{u} = 0$ zorlanır.
\end{enumerate}

\subsection{Parametre Bütçesi}

\begin{table}[H]
\centering
\caption{Yol C parametre tahmini (kuadratik versiyon, 4 blok $\times$ 5 katman = 20 katman).}
\label{tab:yol_c_params}
\begin{tabular}{@{}lc@{}}
\toprule
\textbf{Bileşen} & \textbf{Parametre Sayısı} \\
\midrule
5 operatör $\times$ 4 bileşen $(u,v,w,T')$ $\times$ 13 param $\times$ 20 katman & 5,200 \\
EddyViscosity $C_s$ & 1 \\
$\Delta t$ per katman & 20 \\
\midrule
\textbf{Toplam} & $\sim$5,221 \\
\bottomrule
\end{tabular}
\end{table}

Bu $\sim$5K parametre, bir FNO'nun $\sim$10$^6$ parametresinin $\sim$200 katı altındadır; ancak her parametrenin fiziksel anlamı vardır ve lokal adaptasyon yoluyla $983{,}040$ grid noktasında farklı etkin değerler üretir.


% ============================================================
\section{Bölüm Özeti}
\label{sec:mimari_ozet}

\begin{itemize}
    \item INNATE, NS denklemlerinin her terimini bir nöron olarak kodlar.
    \item 8 nöron tipi: Advection3D, Diffusion, EddyViscosity3D, Projection3D, Forcing3D, Buoyancy3D, ThermalAdvection3D, ThermalDiffusion3D.
    \item Her nöronün 0--2 öğrenilebilir parametresi var; hepsi fiziksel anlamlaştırılabilir.
    \item SpectralOps3DAniso, anizotropik domain için türev/Poisson/dealiasing sağlar.
    \item Forward pass: RHS hesabı $\to$ RK4 integrasyon $\to$ Projection $\to$ gauge fix.
    \item Yol C ile parametreler lokal invariantlara bağlanarak uzamsal adaptasyon elde edilir.
    \item Toplam $\sim$5--10K parametre, tamamı yorumlanabilir.
\end{itemize}

Sonraki bölümde (Bölüm~\ref{ch:training}), bu parametrelerin nasıl öğrenildiğini --- loss fonksiyonları, curriculum learning ve eğitim stratejisini --- inceleyeceğiz.
